% Folder 83, batch 01 (archivists' pages 1-20).
% Pass: Opus 5.5 (claude-opus-5-5), 2026-10-10 — first pass, unchecked against the pages by a human.
%
% Page 1 is blank. Pages 4, 6, 8, 10, 12, 14 and 16 are typescript leaves of
% his non-mathematical writing (Récoltes et Semailles « L » letter pages, and
% « PU » pages on dreams), annotated in his hand; not the mathematics in hand,
% so skipped. The lower half of page 2 is such a typescript, upside down,
% and is left out. Page 17 is an orange folder cover carrying only a pencil
% « sl(2, Λ) »; skipped, recorded once in a note at page 18.
% Pages 2-15: the icosahedron over F_4 (alpha^2 + alpha + 1 = 0, alpha' the
% conjugate root), S = the twelve u in M = End(V) with Tr u = 1, det u = alpha.
% Pages 18-20 (brown ink, dated Nov. 1982 and 19.6.83): the braid relation
% e0 e1 e0 = e1 e0 e1, sigma, rho_i, rho'_i, and SL(2,Z), GL(2,Z).
\documentclass[11pt,a4paper]{article}
\input{../preamble/grothendieck}

\folder{83}
\batch{1}
\pages{1}{20}
\foldertitle{[Icosaèdres, bi-icosaèdres et algèbres d'Azumaya de rang 4] : notes manuscrites (1982-1983, 1986, s.d.).}
\dating{1982-1986}
\watermark{Édition de démonstration}

\begin{document}

\section{L'icosaèdre sur $\mathbb{F}_4$}

\page{2}
\note{feuillet de calculs sans texte suivi. Dans tout ce qui suit jusqu'à la p.~15, $\alpha$ est une racine de $\alpha^2+\alpha+1=0$, $\alpha'$ l'autre, et l'on calcule en caractéristique $2$ ($-1=1$) ; $u'$ désigne le conjugué $\operatorname{Tr}u - u$.}

\marginal{Icosaèdre sur $\mathbb{F}_4$ — \uncertain{gribouillis}}

\[
\alpha^2 \text{\struck{$+$}}\, \alpha + 1 = 0 \qquad (\text{racines } \alpha,\ \alpha')
\]
\[
u^2 = u + \alpha
\]
\[
u^4 = u^2 + \alpha^2 = u^2 + \alpha' = u + \alpha + \alpha' = u + 1 = u'
\]
\[
u^5 = u u' = \alpha\, 1
\]

$u$ d'ordre $15$.

\struck{\ill{}} \quad $u \,/\, u'$

\[
u^3 = u^2 + \alpha u = u + \alpha + \alpha u = \alpha' u + \alpha
\]

$u^3$ d'ordre cinq \quad $\boxed{\alpha' u + \alpha}$ \struck{\ill{}} \quad \uncertain{opérant} \ill{} sur $S \smallsetminus \{u, u'\}$ avec 2 orbites

$u^5$ d'ordre 3

\struck{\ill{}} \quad \uncertain{On veut} $u$ \ill{}

\[
\det(u-v) = 0 \quad \text{i.e.} \quad
\underbrace{\det u}_{1} + \underbrace{\det v}_{1} + \operatorname{Tr} u v = 0
\]

i.e. $\operatorname{Tr} u v = 0$ \qquad $x \otimes x$

N.B.
\[
\operatorname{Tr} u u' = 0
\]
\[
\operatorname{Tr} u.u = \operatorname{Tr} u^2 = \operatorname{Tr}(u + \alpha.1)
= \operatorname{Tr} u = 1 \neq 0
\]
\[
\det(u+v) = \det u + \det v + \operatorname{Tr} u v = \operatorname{Tr} u v
\qquad \text{\struck{$\det(u-v)$}}
\]
\[
\operatorname{Tr} u v = 0
\]
\[
\operatorname{Tr} u v = 0 \qquad
\operatorname{Tr}(u v') = \operatorname{Tr}(u(v+1))
= \operatorname{Tr} u v + \operatorname{Tr} u = 1 + \operatorname{Tr} u v
\]

$\operatorname{Tr} u v$ et $\operatorname{Tr} u v'$ sont \emph{différents}, donc
$\operatorname{Tr} u v = 0 \Rightarrow \operatorname{Tr} u v' \neq 0$.

\[
\operatorname{Tr} u u' = 0, \qquad \operatorname{Tr} u^2 = 1
\]

\page{3}
\note{feuillet de calculs épars ; l'ordre de lecture adopté est celui du haut vers le bas.}

\[
\begin{pmatrix} a & b \\ c & d \end{pmatrix}
\quad \text{\struck{\ill{}}} \quad
\begin{pmatrix} d & -b \\ -c & a \end{pmatrix} \frac{1}{ad-bc}
\]
\[
\begin{pmatrix} a & b \\ c & d \end{pmatrix}
\begin{pmatrix} a' & b' \\ c' & d' \end{pmatrix}
= \begin{pmatrix} aa' + bc' & ab' + bd' \\ \cdots & \cdots \end{pmatrix}
\]
\note{seule la première ligne du produit est écrite.}

\[
\boxed{\begin{aligned}
u u' &= u' u = \det u \\
u + u' &= \operatorname{Tr} u . \mathrm{id}
\end{aligned}}
\qquad 2 \operatorname{Tr} u = 2R
\]
\note{le second membre de « $2\operatorname{Tr}u = 2\ldots$ » est d'une lecture incertaine.}

\struck{$2\operatorname{Tr}u =$} \quad si $\operatorname{Tr} u = 0$ alors

\[
u' = -u \iff \operatorname{Tr} u = 0
\]

\[
u^2 - (\operatorname{Tr} u)\, u + \det u = 0
\]
\[
\operatorname{Tr} u = 1 \qquad \det u = 0 \qquad \det u \overset{?}{=} 1
\qquad u = u'
\]
\[
1 + a' \otimes a, \qquad \langle a, a' \rangle = 1
\]
\[
u \longmapsto -u' \qquad u' u = 0
\]
\[
\text{\struck{$u(x) \wedge y = x \wedge \ldots$}}
\]
\[
u(x) \wedge y = -u'(y) \wedge x
\]
\[
\text{\struck{$\boxed{u(x) \wedge y + x \wedge u'(y) = 0}$}}
\]
\[
\boxed{u(x) \wedge y = x \wedge u'(y)}
\]
\[
E \to \check{E} \simeq E \otimes \omega^{-1}, \qquad
E \times E \to \omega, \qquad
E \simeq \check{E} \otimes \omega, \qquad \check{E}
\]
\[
\lambda e + \mu f, \qquad x \wedge x, \qquad
\text{\struck{$x' \otimes x$}}
\]
\[
15 \cdot 4 = 60, \qquad 4 \cdot 5 = 20
\]

\emph{les projections}

\page{5}
\[
\operatorname{Tr} u = 1 \quad
\begin{array}{|lll r}
\det u = 0 & u^2 - u = 0 & u \text{ un projecteur} & 20 \\
\det u = 1 & u^2 - u + 1 = 0 & u \text{ val. propres } \alpha, \alpha' & 20 \\
\det u = \alpha & u^2 - u + \alpha = 0 & & 12 \\
\det u = \alpha' & u^2 - u + \alpha' = 0 & & 12 \\
 & & & 64
\end{array}
\]
\note{à la ligne $\det u = 1$, « $u = \mathrm{id} + \text{nilpotent}$ » est biffé et remplacé par « $u$ \uncertain{val. propres} $\alpha, \alpha'$ » ; sous $u^2-u+1=0$ est amorcé « $(u+\alpha)(u+\ldots$ ». Les deux « 12 » sont cerclés.}

N.B. Les $u \sim \alpha$ d'ordre 15.

\[
u^2 \text{\struck{$+$}}\, \lambda u + \mu = 0
\]
\[
1,\ u \qquad \alpha.1 \longmapsto \alpha 1 + u \qquad
u \longmapsto \mu.1 + \lambda u
\]
\[
\begin{pmatrix} 0 & -\mu \\ 1 & \lambda \end{pmatrix}
\qquad 1,\ u,\ u'
\]

\[
\operatorname{Tr} u = 1, \quad \operatorname{Tr} v = 1, \qquad
\det u = \det v = \alpha
\]
\[
\operatorname{Tr}(u+v) = 0
\]
\[
\det(u+v) = \underbrace{\det u}_{1} + \underbrace{\det v}_{1}
+ \operatorname{Tr} u v = \operatorname{Tr}(uv)
\]
\note{les accolades « 1 » sous $\det u$, $\det v$ sont de la page ; elles ne s'accordent pas avec $\det u = \det v = \alpha$ écrit juste au-dessus.}

\[
\operatorname{Tr} u v = \left\{
\begin{array}{l} 0 \\ 1 \\ \alpha \\ \alpha' \end{array}
\right.
\qquad
\begin{array}{|l}
\lambda u + \mu \\
\lambda u' + \mu \\
\lambda^2 \alpha + \mu^2 + \lambda\mu \neq 0
\end{array}
\]

\drawing{un carré encadré : $u$ — $uv$ en haut, $1$ — $v$ en bas, avec deux diagonales parallèles $1$–$u$ et $v$–$uv$.}

\[
u = u' \Longrightarrow \quad u \quad u' \quad 1
\]

\[
\lambda u + \mu \qquad
\operatorname{Tr}(\lambda u + \mu) = \lambda \operatorname{Tr} u = \lambda = 1
\]
\[
u + \mu \qquad
\det(u+\mu) = \det u + \mu^2 + \mu \operatorname{Tr} u
= \mu^2 + \mu + \alpha = \text{\struck{$\alpha$}}
\]
\[
u + u' = 1 \qquad u + 1 = u'
\]
\[
\text{\struck{$\mu^2 + \alpha\mu \ldots$}} \qquad
\text{\struck{$\alpha^2 = \alpha + \alpha\alpha'$}} \qquad
\alpha'^2 + \alpha\alpha' + 1 = 0
\]
\[
\mu^2 + \mu + \alpha = \alpha, \qquad \mu^2 + \mu = 0
\]
\[
\operatorname{Tr}(uv) = \qquad u \longmapsto u' = u + 1
\qquad u \longmapsto u' \ \text{antipodisme}
\]
\[
uv = \text{\struck{$u+$}}\ \lambda u + \mu, \qquad
u(v - \lambda) = \mu
\]
\[
v - \lambda = \mu u^{-1} = (\mu\alpha')(u+1), \qquad
v = \mu\alpha' u + \ldots
\]
\[
u' u = \alpha. 1, \qquad u^{-1} = \alpha' u'
\]
\[
\operatorname{Tr} u u' = 0, \qquad \operatorname{Tr} u w = 0, \qquad
v = u' + w, \qquad \operatorname{Tr} w = 0, \qquad
\operatorname{Tr} \text{\struck{$u$}}\, w = 0
\]

\page{7}
\[
\operatorname{Tr} \text{\struck{$v$}}\, u v u^{-1} = 0 \text{ ou } 1
\]
\[
\operatorname{Tr} v\, u(v) = \det(v - u(v)) \in \{0, 1\}
\]
\note{ici et plus loin $u(v)$ désigne le conjugué $u v u^{-1}$.}

$\operatorname{Tr}(uv) \in \{0,1\}$ \quad $\forall u, v \in S$ \qquad
N.B. $\operatorname{Tr} u^2 = 1$, $\operatorname{Tr} u u' = 0$.

\uncertain{Oh oui}, $v = u$ ou $v = u'$ \uncertain{d'accord} \ill{},
$\operatorname{Tr} u v' = \operatorname{Tr} u' v = \operatorname{Tr} u v + 1 \neq \operatorname{Tr} u v$.
Pour les dix autres, \uncertain{avec} \ill{} \uncertain{formant} \ill{} sous le groupe $k[u]^*/k^*$ opérant \ill{} or $k[u]^*$ est cyclique d'ordre $15$ ($4^2 - 1 = 15$), $k^*$ est son groupe d'ordre $3$, donc $k[u]^*/k^*$ est cyclique d'ordre $5$. \uncertain{On compare} que $u^5 = \alpha.1$ donc $u$ est d'ordre $15$ ($\alpha$ étant d'ordre $3$) donc $u$ \add{définit} un \uncertain{générateur} de $k[u]^*/k^*$. Ceci dit \struck{\ill{}} \add{\ill{}} \uncertain{les points fixes dans} $\mathrm{End}(V) = M$ \struck{\ill{}} des $k[u]^*/k^*$, i.e. les commutants de $k[u]$ dans $M$, \ill{} de $u$ est \uncertain{réduit à} $k[u]$. On a vu que
\[
k[u] \cap S = \{u, u'\}
\]
donc $u$ ne fixe aucun pt de $S \smallsetminus \{u, u'\}$, qui se décompose donc en deux orbites d'ordre $5$. Donc $\operatorname{Tr}(uv)$ pour $v \in S \smallsetminus \{u, u'\}$ ne peut prendre que deux valeurs au plus (une par orbite).

\marginal{Cor. 1. $M^*/k^*$ \uncertain{est transitif} \ill{} — Cor. 2. \ill{} sur $S$ \ill{} \uncertain{transitif}}

\page{9}
\uncertain{Soit} l'un des $(u, v) \in S \times S$ tels que $v \neq u$, $\operatorname{Tr} u v = \beta \in I$ \struck{\ill{}} \add{là où} \ill{} $\mathbb{P}_2(k)$ \struck{\ill{}} \ill{} \uncertain{famille} des $\operatorname{Tr} u v$ \quad $u, v \in S$, $v \neq u, u'$ \quad \ill{} \uncertain{donné} \ill{} dans
$\{0, 1\} = \mathbb{F}_2 \subset k$.

Je dis que $\beta \in$ \ldots
\note{la phrase s'arrête là.}

\[
u = \begin{pmatrix} 0 & 1 \\ \alpha & 1 \end{pmatrix}
\qquad
u' = \begin{pmatrix} 1 & 1 \\ \alpha & 0 \end{pmatrix}
\qquad
\begin{pmatrix} a & b \\ c & d \end{pmatrix}
\begin{pmatrix} a' & b' \\ c' & d' \end{pmatrix}
\]
\note{dans $u$ et $u'$, le « 1 » de la première ligne est écrit sur un « $-1$ ».}
\[
u^{-1} = \alpha' \begin{pmatrix} 1 & 1 \\ \alpha & 0 \end{pmatrix}
= \begin{pmatrix} \alpha' & \alpha' \\ 1 & 0 \end{pmatrix}
\qquad e_1, \ e_2 \qquad \operatorname{Tr} u
\]
\[
\begin{pmatrix} \lambda & 0 \\ 0 & 1 \end{pmatrix}
\begin{pmatrix} 0 & 1 \\ \alpha & 1 \end{pmatrix}
\begin{pmatrix} \lambda^{-1} & 0 \\ 0 & 1 \end{pmatrix}
=
\begin{pmatrix} 0 & \lambda \\ \lambda^{-1}\alpha & 1 \end{pmatrix}
\qquad
v_0 = \begin{pmatrix} 0 & -\lambda \\ \lambda^{-1}\alpha & 1 \end{pmatrix}
\]

$\lambda \neq 1$ \quad i.e. $\lambda \in \{\alpha, \alpha'\}$

\[
v \quad u v u^{-1}, \qquad v' \quad u v' u^{-1}, \qquad
u' u = \det u = \alpha
\]
\[
v u v u^{-1} = v u v (\alpha' u'), \qquad
u(u+1) = \alpha, \qquad u^{-1} = \alpha'(u+1),
\qquad (vu)(vu) + vu
\]
\[
v u = \begin{pmatrix} \lambda\alpha & \lambda \\ \alpha & 1 + \lambda^{-1}\alpha \end{pmatrix}
\qquad
v u^{-1} = \begin{pmatrix} \lambda & 0 \\ 1 + \lambda^{-1} & \lambda^{-1} \end{pmatrix}
\]
\note{dans $vu^{-1}$, le coefficient $\lambda$ est écrit sur un « $\alpha\lambda$ » biffé.}
\[
\lambda^2 \alpha + \lambda + 1 + \lambda^{-1} + \lambda^{-2}\alpha
= \alpha(\lambda^2 + \lambda^{-2}) + (\lambda + \lambda^{-1}) + 1
\]
\[
\text{\struck{$\lambda\alpha +$}} \quad \lambda + \lambda^{-1} \overset{?}{=} 1
\qquad \boxed{\alpha\sigma^2 + \sigma + 1} \qquad \alpha .
\]
\note{$\sigma$ abrège ici, sans être défini sur la page, $\lambda + \lambda^{-1}$ (on a $\lambda^2 + \lambda^{-2} = \sigma^2$ en caractéristique $2$).}

\page{11}
\[
v = \lambda u + \mu \qquad \operatorname{Tr} v = 1 \quad \text{i.e. } \lambda = 1
\]
\[
v = u + \mu, \qquad
\det(v) = \underbrace{\det u}_{1} + \det \mu.1 + \operatorname{Tr} \mu u
= \alpha + \mu^2 + \mu
\]

La condition $\det(v) = \alpha$ équivaut à
\[
\mu^2 + \mu = 0 \qquad \mu \in \{0, 1\} \quad \text{i.e.} \quad v \in \{u, u'\}
\]
La condition $\det(v) = \alpha'$ équivaut à
\[
\mu^2 + \mu + 1 = 0 \qquad \text{i.e.} \quad \mu \in \{\alpha, \alpha'\}
\]
\[
v = u + \alpha, \qquad \underbrace{u + \alpha'}_{u' + \alpha} \in S'
\qquad 0 \quad 1 \quad \alpha \quad \alpha' \ \ldots \quad \alpha^2 = \alpha'
\]
\[
\operatorname{Tr} u v = 0
\]
\[
\operatorname{Tr}(u + \alpha)(v + \alpha) = \operatorname{Tr} u v
+ \text{\struck{\ill{}}}
\qquad uv + \alpha u + \alpha v + \alpha^2
\]
\[
\operatorname{Tr}(u + \alpha, v) = \operatorname{Tr} u v + \alpha
\qquad v \quad u v u^{-1}
\]
\[
\text{\struck{$\operatorname{Tr} v (u v u^{-1}) = \operatorname{Tr}(v u)^2 u^{-2} v'$}}
\qquad u(v)' = u(v')
\]

\[
\begin{array}{c|cccc}
\lambda & 0 & 1 & \alpha & \alpha' \\ \hline
\sigma & & & &
\end{array}
\]

\drawing{un triangle ; puis un pentagone de sommets reliés à un centre (l'étoile d'un sommet de l'icosaèdre), une face marquée $\alpha$.}

\[
\boxed{\operatorname{Tr} v_0\, u(v_0) = \alpha}
\Longrightarrow
\operatorname{Tr} v\, u(v) = \alpha
\]
pour les \uncertain{cinq conjugués} de $v$.
\[
\operatorname{Tr} v'\, \underbrace{u(v')}_{u(v)'} = \alpha
\]
\[
\det(u + v) = \operatorname{Tr}(uv) =
\left\{
\begin{array}{ll}
\alpha & \text{si } \{u, v\} \text{ arête} \\
\alpha' & \text{si } \{u, v\} \text{ coarête} \\
1 & \text{si } u = v \\
0 & \text{si } u = u'
\end{array}
\right.
\]
\note{les mots qualifiant la paire $\{u,v\}$ sont de lecture douteuse ; la valeur « 1 » est écrite sur un chiffre biffé, le « 0 » aussi.}

\page{13}
\note{changement de cadre : la page calcule avec des signes, en caractéristique quelconque.}

\drawing{dans le plan, les vecteurs $e_1$, $e_2$, $e_3 = e_1 + e_2$ (diagonale du carré) et $e_4 = e_1 + \lambda e_2$.}

\[
e_3 = \begin{pmatrix} 1 \\ 1 \end{pmatrix}, \qquad
e_4 = \begin{pmatrix} 1 \\ \lambda \end{pmatrix}
\]
\[
\boxed{u = \begin{pmatrix} \zeta & 0 \\ 0 & -\zeta \end{pmatrix}}
\qquad \operatorname{Tr} u = 0, \quad \det u = -\zeta^2;
\qquad \operatorname{Tr} v = 0, \quad \det v = -\eta^2
\]
\[
v = \begin{pmatrix} a & b \\ c & d \end{pmatrix}, \qquad
v(e_1) = a e_1 + c e_2
\]
\[
v(e_3) = (a + b,\ c + d) = (\eta, \eta), \qquad
v(e_4) = (a + \lambda b,\ c + \lambda d) = (-\eta, -\lambda\eta)
\]
\[
\begin{aligned}
a + b &= \eta, & c + d &= \eta, \\
a + \lambda b &= -\eta, & c + \lambda d &= -\lambda\eta
\end{aligned}
\]
\[
b = \eta - a, \quad d = \eta - c
\]
\note{à gauche de ces deux égalités, une matrice $v = \left(\begin{smallmatrix} a & c \\ c & a \end{smallmatrix}\right)$ est biffée.}
\[
a + \lambda(\eta - a) = -\eta, \quad (1 - \lambda) a = -(1 + \lambda)\eta;
\qquad
c + \lambda(\eta - c) = -\lambda\eta, \quad (1 - \lambda) c = -2\lambda\eta
\]
\[
a = \frac{\lambda + 1}{\lambda - 1}\,\eta, \qquad
c = \frac{2\lambda\eta}{\lambda - 1}
\]
\[
b = \eta\Bigl(1 - \frac{\lambda + 1}{\lambda - 1}\Bigr) = \frac{-2\eta}{\lambda - 1},
\qquad
d = \eta\Bigl(1 - \frac{2\lambda}{\lambda - 1}\Bigr)
= -\frac{\lambda + 1}{\lambda - 1}\,\eta = -a,
\qquad c = -\lambda b
\]
\[
\boxed{v = \begin{pmatrix} a & b \\ -\lambda b & -a \end{pmatrix}}
\qquad
\left\{
\begin{aligned}
a &= \frac{\lambda + 1}{\lambda - 1}\,\eta \\
b &= \frac{-2}{\lambda - 1}\,\eta
\end{aligned}
\right.
\]
\[
\det v = -a^2 + \lambda b^2
= -\eta^2 \underbrace{\Bigl[\Bigl(\frac{\lambda + 1}{\lambda - 1}\Bigr)^2
- \lambda \frac{4}{(\lambda - 1)^2}\Bigr]}_{
\frac{(\lambda + 1)^2 - 4\lambda}{(\lambda - 1)^2} = 1}
\]
\[
\text{\struck{$uv + vu$}} \qquad
u v = \begin{pmatrix} \zeta a & \zeta b \\ \lambda\zeta b & \zeta a \end{pmatrix}
\qquad
v u = \begin{pmatrix} \zeta a & -\zeta b \\ -\lambda\zeta b & \zeta a \end{pmatrix}
\]
\note{le signe devant $\lambda\zeta b$ dans $uv$ est mal formé.}
\[
uv + vu = \begin{pmatrix} 2\zeta a & 0 \\ 0 & 2\zeta a \end{pmatrix}
= \text{\struck{$2\zeta a \ldots$}}\ 2\zeta a . 1
\qquad 2\zeta a = \operatorname{Tr} u v
\]

\page{15}
\[
\begin{array}{c|c|c|c|c}
(u, v) & v = u & v = u' & \{u, v\} \in A_\alpha & \{u, v\} \in A'_\alpha \\ \hline
\det(u - v) & 0 & 1 & \alpha' & \alpha \\ \hline
\operatorname{Tr} u v & 1 & 0 & \alpha & \alpha'
\end{array}
\]

\textbf{Prop.} Soient $u, v \in S_\alpha$, avec $v \neq u, u'$. Alors
\[
\begin{aligned}
\operatorname{Tr} v\, u^i(v) &= \alpha & &\text{si } i \equiv 1 \text{ ou } -1 \ (5) \\
\operatorname{Tr} v\, u^i(v) &= \alpha' & &\text{si } i \equiv 2 \text{ ou } -2 \ (5) \\
(&= 1 & &\text{si } i \equiv 0 \ (5))
\end{aligned}
\]

\drawing{à gauche de l'énoncé, l'étoile pentagonale d'un sommet $u_0$, avec une arête épaissie vers $u_1$.}

\[
\operatorname{Tr} u v_0 = 1 + \alpha(\lambda + \lambda^{-1}) = 1 + \alpha = \alpha'
\qquad u^5 = \alpha
\]
\[
\operatorname{Tr} u v_0' = \alpha
\qquad
u v = \begin{pmatrix} 1 & 0 \\ \alpha' & \alpha' \end{pmatrix}
\qquad
u^{-1} v = \begin{pmatrix} 0 & 1 \\ 1 & \alpha' \end{pmatrix}
\]
\[
u = \begin{pmatrix} 0 & 1 \\ \alpha & 1 \end{pmatrix}
\qquad
u^{-1} = \begin{pmatrix} \alpha' & \alpha' \\ 1 & 0 \end{pmatrix}
\qquad
v_0 = \begin{pmatrix} 1 & \alpha \\ 1 & 0 \end{pmatrix}
\]
\[
v_1 = u v u^{-1} = \begin{pmatrix} \alpha' & \alpha \\ 1 & \alpha' \end{pmatrix}
\qquad
v_{-1} = u^{-1} v u = \begin{pmatrix} \alpha & 1 \\ \alpha' & \alpha' \end{pmatrix}
\]
\note{plusieurs coefficients de ces matrices sont récrits sur d'autres ; on donne la dernière lecture.}

\drawing{l'étoile du sommet $u$ : un pentagone $v_0, v_1, \ldots, v_{-1} = v_4$, une flèche de rotation au centre.}

\[
u^5 = \alpha, \qquad (\alpha u)^5 = \alpha^6 = 1, \qquad
\det(\alpha u) = \alpha^3 = 1
\]
\uncertain{v.\,l. propres} de $u$ \ill{} : $\alpha, \alpha'$ \quad \struck{$\alpha u$}

\[
\frac{12 \cdot 11}{2} = 66 = 6 + 30 + 30, \qquad
\frac{12 \cdot 11 \cdot 10}{6} = 220 = 20 + 20 + 60 + 60 + 6 \cdot 10
\qquad 1 + \alpha = \alpha'
\]

\drawing{un réseau de faces de l'icosaèdre autour de $u$ et $u_0$, sommets marqués $v_1$, $v_3'$, $v_4'$, $v_2$, $v_6'$, $v_3$, $v_4$, $v_2'$, $v_1'$ ; une face hachurée.}

\[
\operatorname{Tr} \text{\struck{$v$}}\, u w u^{-1} = \alpha
\qquad \text{\struck{\ill{}}}\ u' \ \text{\struck{\ill{}}}
\]

$(u, v, w)$ formant un triangle,
\[
\det(u + v + w) = \underbrace{\det u + \det v + \det w}_{3\alpha = \alpha}
+ \underbrace{\varphi(u, v) + \varphi(v, w) + \varphi(w, u)}_{3\alpha' = \alpha'}
= 1
\]

\section{La relation $\varepsilon_0\varepsilon_1\varepsilon_0 = \varepsilon_1\varepsilon_0\varepsilon_1$}

\page{18}
\note{la chemise orange qui précède (p.~17, non transcrite) porte seulement, au crayon, « $\mathrm{sl}(2, \Lambda)$ », d'une main qu'on ne peut attribuer avec certitude. Les pp.~18–20 sont à l'encre brune.}

\marginal{Nov. 1982}

$\varepsilon_0, \varepsilon_1 \in H$ un groupe. OPS $\varepsilon_0, \varepsilon_1$ engendrant.

posons
\[
\left\{
\begin{aligned}
\sigma &= \varepsilon_0 \varepsilon_1 \varepsilon_0 \quad (\text{dissymétrique}) \\
\rho_0' &= \sigma \varepsilon_0, \quad \rho_1' = \sigma \varepsilon_1 \\
\rho_0 &= \sigma^{-1} \varepsilon_0 = \varepsilon_0^{-1}\varepsilon_1^{-1}
= (\varepsilon_1 \varepsilon_0)^{-1}, \quad
\rho_1 = \varepsilon_1^{-1}\varepsilon_0^{-1}
\end{aligned}
\right.
\]

on a alors \uncertain{trivialement}
\[
\left\{
\begin{aligned}
&\varepsilon_0 \rho_0 = \rho_1 \varepsilon_0 = \rho_0' \varepsilon_0^{-1}
= \rho_1' \varepsilon_1^{-1}, \quad \text{donc } \rho_1 = \varepsilon_0(\rho_0) \\
&\rho_1' = \rho_1^{-2} \\
&\sigma(\varepsilon_0) = \rho_0'(\varepsilon_0)
\end{aligned}
\right.
\]
\note{la fin de la dernière ligne (« $= \rho_1(\varepsilon_0) = \ldots$ ») et une ligne au-dessous sont biffées ; à droite, le calcul « $\varepsilon_0\varepsilon_1\varepsilon_0\varepsilon_1^{-1}\varepsilon_0^{-1} = \varepsilon_1$ » est suivi de « \uncertain{Où si l'on suppose} $\varepsilon_0\varepsilon_1\varepsilon_0 = \varepsilon_1\varepsilon_0\varepsilon_1$ ».}

\note{la suite de la page est une liste de formules étiquetées a) à j), la colonne de gauche disant ce que chaque groupe exprime ; les étiquettes c) à g) sont récrites sur d'autres, décalées d'une lettre.}

$(\rho_0, \rho_0')$ \quad $(\rho_1, \rho_1')$ \quad $\sigma$ \quad $\varepsilon_0$ \quad $\varepsilon_1$

a) \emph{symétrie : conj. par $\sigma$}
\[
\sigma(\varepsilon_0) = \varepsilon_1, \quad \sigma(\varepsilon_1) = \varepsilon_0;
\qquad
\sigma(\rho_0) = \rho_1, \quad \sigma(\rho_1) = \rho_0, \quad
\sigma(\rho_0') = \rho_1', \quad \sigma(\rho_1') = \rho_0'
\]

a') \emph{conjug.}
\[
\begin{aligned}
&\rho_0(\varepsilon_0) = \rho_0'(\varepsilon_0) = \varepsilon_1, \qquad
\rho_1(\varepsilon_1) = \rho_1'(\varepsilon_1) = \varepsilon_0, \\
&\varepsilon_0(\rho_0) = \rho_1, \quad \varepsilon_0(\rho_0') = \rho_1', \qquad
\varepsilon_1(\rho_1) = \rho_0, \quad \varepsilon_1(\rho_1') = \rho_0'
\end{aligned}
\]

\marginal{$\sigma^2$ commute à $\varepsilon_0, \varepsilon_1, \rho_0, \rho_0', \rho_1, \rho_1'$, soit $\sigma^2 = \omega$ \quad (au-dessus, biffé : $\sigma^2(\varepsilon_0) = \varepsilon_0$, $\sigma^2(\varepsilon_1) = \varepsilon_1$)}

b) \emph{introduction des $\rho_i$ en termes des $\sigma$, $\varepsilon_i$}
\[
\rho_0 = \sigma^{-1}\varepsilon_0 = \varepsilon_1 \sigma^{-1}, \quad
\rho_1 = \sigma^{-1}\varepsilon_1 = \varepsilon_0 \sigma^{-1};
\qquad
\rho_0' = \sigma \varepsilon_0 = \varepsilon_1 \sigma, \quad
\rho_1' = \sigma \varepsilon_1 = \varepsilon_0 \sigma
\]

c) \emph{éléments d'ordre fini}
\[
\rho_0^3 = \rho_1^3 = \sigma^{-2} = \omega, \qquad \omega^2 = 1
\]

d)
\[
\Bigl[\ \rho_0^6 = \rho_1^6 = \sigma^{-4} = 1, \qquad
\rho_0'^3 = \rho_1'^3 = 1 \ \Bigr]
\]
\marginal{$\rho_0' = \rho_0^{-2}$ \ldots\ valables dans $\mathrm{Sl}(2, \mathbb{Z})$ mais non \ill{} $\mathrm{Sl}(2, \mathbb{Z})$}

e) \emph{relations entre $\rho_i$, $\rho_i'$}
\[
\rho_0' = \omega \rho_0, \quad \rho_1' = \omega \rho_1; \qquad
\rho_0' = \rho_0^{-2}, \quad \rho_1' = \rho_1^{-2}
\qquad \text{d'où } \rho_0 = \omega^{-1}\rho_0', \ \rho_1 = \omega^{-1}\rho_1'
\]
\note{une ligne entre les deux est biffée et illisible.}

f) \emph{expressions en $\varepsilon_0$, $\varepsilon_1$}
\[
\begin{aligned}
\sigma &= \varepsilon_0 \varepsilon_1 \varepsilon_0
\ [= \varepsilon_1 \varepsilon_0 \varepsilon_1] \\
\rho_0 &= \varepsilon_0^{-1}\varepsilon_1^{-1} = (\varepsilon_1 \varepsilon_0)^{-1},
& \rho_1 &= \varepsilon_1^{-1}\varepsilon_0^{-1} = (\varepsilon_0 \varepsilon_1)^{-1} \\
\rho_0^{-1} &= \varepsilon_1 \varepsilon_0, & \rho_1^{-1} &= \varepsilon_0 \varepsilon_1
\end{aligned}
\]

g) \emph{expressions en $\sigma$, $\varepsilon_0$}
\[
\begin{aligned}
\varepsilon_1 &= \sigma(\varepsilon_0) = \sigma \varepsilon_0 \sigma^{-1} \\
\rho_0 &= \sigma^{-1}\varepsilon_0, & \rho_0' &= \sigma \varepsilon_0 \\
\rho_1 &= \varepsilon_0 \sigma^{-1}, & \rho_1' &= \varepsilon_0 \sigma
\end{aligned}
\]

h) \emph{expressions en $(\sigma, \rho_0)$}
\[
\begin{aligned}
\varepsilon_0 &= \sigma \rho_0, \qquad \varepsilon_1 = \rho_0 \sigma, \qquad
\rho_0' = \rho_0^{-2} \\
\rho_1 &= \sigma(\rho_0) = \sigma \rho_0 \sigma^{-1}, \qquad
\rho_1' = \sigma(\rho_0^{-2}) = \sigma \rho_0^{-2} \sigma^{-1}
\end{aligned}
\]

i) \emph{expressions en $(\sigma, \rho_0')$}
\[
\begin{aligned}
\varepsilon_0 &= \sigma^{-1}\rho_0', \qquad \varepsilon_1 = \rho_0' \sigma^{-1}, \qquad
\rho_0 = \sigma^{-2}\rho_0' = \rho_0' \sigma^{-2} \\
\rho_1 &= \sigma^{-1}\rho_0' \sigma^{-1}, \qquad
\rho_1' = \sigma(\rho_0') = \sigma \rho_0' \sigma^{-1}
\end{aligned}
\]

j) \emph{expressions en $(\varepsilon_0, \rho_0)$}
\[
\begin{aligned}
\varepsilon_1 &= \rho_0(\varepsilon_0) = \rho_0 \varepsilon_0 \rho_0^{-1}, \qquad
\sigma = \varepsilon_0 \rho_0^{-1}, \qquad \rho_0' = \rho_0^{-2} \\
\rho_1 &= \varepsilon_0(\rho_0) = \varepsilon_0 \rho_0 \varepsilon_0^{-1}, \qquad
\rho_1' = \varepsilon_0(\rho_0^{-2}) = \varepsilon_0 \rho_0^{-2} \varepsilon_0^{-1}
\end{aligned}
\]

\page{19}
k) \emph{expressions en $(\varepsilon_0, \rho_0')$}
\[
\begin{aligned}
\varepsilon_1 &= \rho_0'(\varepsilon_0) = \rho_0' \varepsilon_0 \rho_0'^{-1}, \qquad
\sigma = \rho_0' \varepsilon_0^{-1}, \qquad
\rho_0 = \varepsilon_0 \rho_0'^{-1} \varepsilon_0 \\
\rho_1 &= \varepsilon_0^2 \rho_0'^{-1}, \qquad
\rho_1' = \varepsilon_0(\rho_0') = \varepsilon_0 \rho_0' \varepsilon_0^{-1}
\end{aligned}
\]

l) \emph{expressions en $(\varepsilon_0, \rho_1)$}
\[
\begin{aligned}
\varepsilon_1 &= \rho_1^{-1}(\varepsilon_0) = \rho_1^{-1} \varepsilon_0 \rho_1, \qquad
\sigma = \rho_1^{-1} \varepsilon_0, \qquad
\rho_0 = \varepsilon_0^{-1}(\rho_1) = \varepsilon_0^{-1} \rho_1 \varepsilon_0 \\
\rho_1' &= \rho_1^{-2}, \qquad
\rho_0' = \varepsilon_0^{-1}(\rho_1^{-2}) = \varepsilon_0^{-1} \rho_1^{-2} \varepsilon_0
\end{aligned}
\]

\uncertain{m)} \emph{expressions en $(\varepsilon_0, \rho_1')$}
\[
\begin{aligned}
\varepsilon_1 &= \rho_1'^{-1}(\varepsilon_0) = \rho_1'^{-1} \varepsilon_0 \rho_1', \qquad
\sigma = \varepsilon_0^{-1} \rho_1', \qquad
\rho_1 = \varepsilon_0 \rho_1'^{-1} \varepsilon_0 \\
\rho_0' &= \varepsilon_0^{-1}(\rho_1') = \varepsilon_0^{-1} \rho_1' \varepsilon_0, \qquad
\rho_0 = \rho_1'^{-1} \varepsilon_0^2
\end{aligned}
\]

\marginal{19.6.83}

N.B. parmi les sept éléments $\rho_i, \rho_i', \varepsilon_i, \sigma$, \struck{deux quelconques} \add{une} paire quelconque qui n'est pas formée uniquement d'éléments parmi les 4 $\rho$ [ce qui fait $\binom{7}{2} - \binom{4}{2} = 21 - 6 = 15$ paires] engendrent $H$. D'autre part, les quatre éléments $\rho$ ($\rho_0, \rho_1, \rho_0', \rho_1'$) \uncertain{n'engendrent} \ill{} pas $H = \langle \varepsilon_0, \varepsilon_1 \mid \varepsilon_0\varepsilon_1\varepsilon_0 = \varepsilon_1\varepsilon_0\varepsilon_1 \rangle$, \ill{} \uncertain{commutant} \ill{} \ldots

Il faudrait aussi \uncertain{expliciter} l'extension canonique de $\mathbb{Z}/2\mathbb{Z}$ par
\[
\begin{aligned}
S\mathcal{T}^{+}_{1,1}
&= \langle \varepsilon_0, \varepsilon_1 \mid
\varepsilon_0 \varepsilon_1 \varepsilon_0 = \varepsilon_1 \varepsilon_0 \varepsilon_1 \rangle \\
&\simeq \langle \rho, \sigma \mid \rho^{-3} = \sigma^2 \rangle \\
&\simeq \langle \rho, \varepsilon_0 \mid
\rho \varepsilon_0 \rho^{-1} \varepsilon_0 \rho = 1 \rangle
\end{aligned}
\]
\note{le symbole lu $S\mathcal{T}^{+}_{1,1}$ est d'une forme incertaine (un $T$ cursif précédé d'un $S$).}
donnant \uncertain{ainsi} (\uncertain{puis} divisant par la relation $\omega = 1$) : $\mathrm{Gl}(2, \mathbb{Z})$ est ext. de $\mathbb{Z}/2\mathbb{Z}$ par $\mathrm{Sl}(2, \mathbb{Z})$.

\page{20}
\note{le paragraphe encadré en tête de page reprend une discussion qui ne figure pas dans ce lot : $U_0$, $T_H$, $T_0$ et $\zeta$ n'y sont pas définis. Elle vient d'ailleurs dans la chemise.}

\ill{} non l'égalité — celle-ci est \ill{} \ill{}. La situation devient \uncertain{encore} plus jolie si on suppose $\zeta = 1$, puisque alors la donnée supplémentaire devient : la donnée d'un sous-groupe \add{(commutatif)} $T_H$ de $H$, normalisant $U_0$, tel que $T_H \cap U_0 = \{1\}$, et que $\sigma(t) = t^{-1}$ pour $t \in T_H$.
\marginal{(\uncertain{eq})}

\begin{tikzcd}
\widetilde{\mathbb{Z}} \arrow[r, "\varepsilon"] \arrow[d, "u_0"'] & \mathbb{Z} \arrow[d, "u"] \\
T_0 \arrow[r, "c_0"] & T_H
\end{tikzcd}

\note{les deux flèches verticales portent le signe $\wr$ (isomorphismes) ; l'étiquette de la flèche du haut est d'une lecture incertaine.}

\[
\lambda^{-1} \quad \lambda^{-\ldots} \quad 1 \quad \lambda^2 \qquad\qquad
\widetilde{\mathbb{Z}}
\]

\[
\varepsilon_0, \ \varepsilon_1, \ \sigma, \ \rho_0, \ \rho_1, \ \rho_0', \ \rho_1', \ \omega
\]
\[
(*)\left\{
\begin{aligned}
&\rho_0' = \sigma \varepsilon_0, \quad \rho_0 = \sigma^{-1}\varepsilon_0 = (\varepsilon_1 \varepsilon_0)^{-1},
\quad \rho_0' = \rho_0^{-2}, \quad \rho_0'^3 = 1, \quad \rho_0^6 = 1 \\
&\rho_1' = \sigma \varepsilon_1, \quad \rho_1 = \sigma^{-1}\varepsilon_1 = (\varepsilon_0 \varepsilon_1)^{-1},
\quad \rho_1' = \rho_1^{-2}, \quad \rho_1'^3 = 1, \quad \rho_1^6 = 1 \\
&\varepsilon_1 = \sigma(\varepsilon_0) = \rho_0(\varepsilon_0) = \rho_0'(\varepsilon_0)
\quad [= \varepsilon_1(\varepsilon_0)] \\
&\varepsilon_0 = \sigma(\varepsilon_1) = \rho_1(\varepsilon_1) = \rho_1'(\varepsilon_1)
\quad [= \varepsilon_0(\varepsilon_1)] \\
&\sigma = \varepsilon_1 \varepsilon_0 \varepsilon_1 = \varepsilon_0 \varepsilon_1 \varepsilon_0
\end{aligned}
\right.
\]
\note{les termes « $= \varepsilon_1(\varepsilon_0)$ », « $= \varepsilon_0(\varepsilon_1)$ » sont barrés par un long trait de crayon qui traverse la page en diagonale ; on ne sait s'il les annule.}

\[
\begin{aligned}
&\rho_0^3 = \rho_1^3 = \sigma^2 = \omega \\
&\omega^2 = 1, \ \omega \text{ commute à tout} \\
&\rho_0' = \omega \rho_0, \quad \rho_0 = \omega \rho_0' \\
&\rho_1' = \omega \rho_1, \quad \rho_1 = \omega \rho_1'
\end{aligned}
\]

Partons de $\varepsilon_0, \varepsilon_1$ [\uncertain{satisfaisant} à (1) $\boxed{\varepsilon_1 \varepsilon_0 \varepsilon_1 = \varepsilon_0 \varepsilon_1 \varepsilon_0}$]; posons $\sigma = \varepsilon_0 \varepsilon_1 \varepsilon_0$ $[= \varepsilon_1 \varepsilon_0 \varepsilon_1]$, $\rho_0' = \sigma \varepsilon_0$, $\rho_1' = \sigma \varepsilon_1$, $\rho_0^{-1} = \varepsilon_1 \varepsilon_0$, $\rho_1^{-1} = \varepsilon_0 \varepsilon_1$ ;
\struck{je dis qu'on a bien} examinons les formules $(*)$ qui restent à vérifier. On a \add{effectivement}
\[
\underbrace{\sigma(\varepsilon_0) = \rho_0'(\varepsilon_0)}
= \varepsilon_0 \varepsilon_1 \varepsilon_0 \varepsilon_1^{-1} \varepsilon_0^{-1}
\overset{?}{=} \varepsilon_1
\]
c'est-à-dire égal à $\varepsilon_1$ ssi \struck{$\varepsilon_1\varepsilon_0 = \varepsilon_0\varepsilon_1$} on a (1) ;
\[
\sigma(\varepsilon_1) = \rho_1'(\varepsilon_1)
= \varepsilon_0 \varepsilon_1 \varepsilon_0 \varepsilon_1 \varepsilon_0^{-1}
\varepsilon_1^{-1} \varepsilon_0^{-1} \overset{?}{=} \varepsilon_0
\]
égal à $\varepsilon_0$ ssi $\varepsilon_1 \varepsilon_0 \varepsilon_1 = \varepsilon_0 \varepsilon_1 \varepsilon_0$, i.e. (1) ;
\[
\rho_0^2 = \varepsilon_1 \varepsilon_0 \varepsilon_1 \varepsilon_0 \overset{?}{=} \rho_0'
= \varepsilon_0 \varepsilon_1 \varepsilon_0 . \varepsilon_0
\quad \text{ssi on a (1)}
\]
\[
\rho_1^2 = \varepsilon_0 \varepsilon_1 \varepsilon_0 \varepsilon_1 \overset{?}{=} \rho_1'
= \varepsilon_0 \varepsilon_1 \varepsilon_0 \varepsilon_1
\quad \text{toujours vrai}
\]
\note{la page écrit $\rho_0^2$, $\rho_1^2$ là où les définitions qui précèdent ($\rho_0^{-1} = \varepsilon_1\varepsilon_0$, $\rho_0' = \rho_0^{-2}$) donneraient $\rho_0^{-2}$, $\rho_1^{-2}$ ; on transcrit tel quel. Le calcul se poursuit vraisemblablement au-delà du lot.}

\end{document}
